% !TeX root = main.tex

\documentclass[a4paper]{article}

% 中文支持与字体选择
\usepackage[no-math]{fontspec}
\usepackage[fontset=none]{ctex}

% 数学支持
\usepackage[tbtags]{amsmath}
\usepackage{unicode-math}
\usepackage[per-mode = symbol]{siunitx}
\usepackage{algpseudocode}
\usepackage{diffcoeff}

% 绘制三线表与跨页表格
\usepackage{booktabs}
\usepackage{longtable}
\usepackage{tabularx}

% 页边距
\usepackage{geometry}

% 页眉与页脚
\usepackage{fancyhdr}

% 图片支持
\usepackage{graphicx}

% 颜色支持
\usepackage{color}
\usepackage[dvipsnames, svgnames, x11names]{xcolor}

% 超链接与书签支持
\usepackage[hidelinks, bookmarksnumbered, bookmarksopen=true]{hyperref}

% 绘图包
\usepackage{float}
\usepackage{tikz}
\usepackage{pgfplots}
\usepackage{circuitikz}

% 目录与标题
\usepackage{titlesec}
\usepackage{titletoc}
\usepackage{caption}

% 参考文献
\usepackage{gbt7714}

% Misc
\usepackage{import}
\usepackage{verbatim}
\usepackage{ulem}
\usepackage{enumitem}

\setlist{nosep}

\setlist[enumerate]{
  label=（\arabic*）,
  left=2\ccwd,
  itemindent=4.5\ccwd
}

\setmainfont{STIXTwoText}[
  UprightFont = *-Regular ,
  BoldFont = *-Bold ,
  ItalicFont = *-Italic ,
  BoldItalicFont = *-BoldItalic ,
  Extension = .otf ,
  Scale = 1.0]

\setsansfont{texgyreheros}[
  UprightFont = *-regular ,
  BoldFont = *-bold ,
  ItalicFont = *-italic ,
  BoldItalicFont = *-bolditalic ,
  Extension = .otf ,
  Scale = 0.9]

\setmonofont{NewCMMono10-Book.otf}[
  ItalicFont=NewCMMono10-BookItalic.otf,
  BoldFont=NewCMMono10-Bold.otf,
  BoldItalicFont=NewCMMono10-BoldOblique.otf,
  SlantedFont=NewCMMono10-Book.otf,
  SlantedFeatures={FakeSlant=0.25},
  BoldSlantedFont=NewCMMono10-Bold.otf,
  BoldSlantedFeatures={FakeSlant=0.25},
  SmallCapsFeatures={Numbers=OldStyle}]

\setmathfont{STIXTwoMath-Regular.otf}

% 页边距与页眉设置
\geometry{a4paper, left = 3.17cm, right = 3.17cm, top = 2.54cm, bottom = 2.54cm}
\setlength{\headheight}{1.13cm}
\setlength{\headsep}{0.3cm}
\fancyhead[C]{\includegraphics[scale = 0.06]{../assets/sspu_logo_title.eps}}
\fancyfoot[C]{}
\fancyhead[LR]{}
\fancyfoot[LR]{}
\pagestyle{fancy}

% 标题设置
\titleformat{\section}[block]{\centering\zihao{3}\sffamily}{\thesection\quad}{0em}{}[]
\titleformat{\subsection}[block]{\zihao{4}\sffamily}{\thesubsection\quad}{0em}{}[]
\titleformat{\subsubsection}[block]{\zihao{-4}\sffamily}{\thesubsubsection\quad}{0em}{}[]

\titlecontents{section}[0bp]{\zihao{4}}{\thecontentslabel\quad}{}{\titlerule*[0.5pc]{.}\contentspage}
\titlecontents{subsection}[2\ccwd]{\zihao{-4}}{\thecontentslabel\quad}{}{\titlerule*[0.5pc]{.}\contentspage}
\titlecontents{subsubsection}[4\ccwd]{\kaishu\zihao{-4}}{\thecontentslabel\quad}{}{\titlerule*[0.5pc]{.}\contentspage}

\captionsetup{font = sf}

% 行距设置
\linespread{1.5}
\setlength{\parskip}{0em}

% 计数器
\newcounter{code}[section]
\numberwithin{figure}{section}
\numberwithin{table}{section}
\numberwithin{equation}{section}

\renewcommand\thecode{\thesection-\arabic{code}}
\renewcommand\thefigure{\thesection-\arabic{figure}}
\renewcommand\thetable{\thesection-\arabic{table}}
\renewcommand\theequation{\thesection-\arabic{equation}}

% 代码环境配置
\newcommand\inputcode[2]{
  \linespread{1}
  \refstepcounter{code}
  \label{#1}
  \begin{figure}[htbp!]
    \vspace{-1em}
    \rule[0.5ex]{\textwidth}{0.08em}
    \noindent{\sffamily{\bfseries 代码片段\thecode：}#1} \\
    \rule[0.5ex]{\textwidth}{0.05em}
    \vspace{-2em}
    \verbatiminput{#2}
    \vspace{-1em}
    \rule[0.5ex]{\textwidth}{0.08em}
  \end{figure}
  \linespread{1.5}
}

% Misc
\pgfplotsset{compat=1.8}
\newcolumntype{Y}{>{\centering\arraybackslash}X}

\bibliographystyle{gbt7714-numerical}

% 英文默认字体：Times New Roman
\setmainfont{Times New Roman}

% 中文默认字体：宋体，粗体以黑体代替，斜体以楷书代替
\setCJKmainfont[BoldFont=FZHei-B01, ItalicFont=FZNewKai-Z03]{FZNewShuSong-Z10}
\setCJKsansfont[AutoFakeBold,AutoFakeSlant]{FZHei-B01}
\setCJKmonofont[AutoFakeBold,AutoFakeSlant]{FZKeBenFangSongS-R-GB}

\newCJKfontfamily\songti[BoldFont={FZHei-B01},ItalicFont={FZNewKai-Z03},AutoFakeBold]{FZNewShuSong-Z10}
\newCJKfontfamily\heiti[AutoFakeBold,AutoFakeSlant]{FZHei-B01}
\newCJKfontfamily\fangsong[AutoFakeBold,AutoFakeSlant]{FZKeBenFangSongS-R-GB}
\newCJKfontfamily\kaishu[AutoFakeBold,AutoFakeSlant]{FZNewKai-Z03}
